%----------------------------------------------------------------------
% 中文摘要
%----------------------------------------------------------------------

% 把中文摘要寫在裡面
\begin{zhAbstract}

    自回歸大型語言模型在推論過程中通常逐一生成詞元，使生成延遲隨輸出序列長度增加。多詞元預測（Multi-Token Prediction, MTP）透過共享模型主幹與多個預測頭，使模型能在單次前向傳播中同時預測多個未來詞元，因而具備加速解碼的潛力。然而，標準 MTP 主要由各預測頭針對其對應的真實詞元獨立計算交叉熵損失，缺乏不同預測頭之間的一致性約束，使距離當前位置較遠的未來詞元預測面臨較高的最佳化難度。

    為改善上述問題，本研究提出選擇性分布一致性（Selective Distribution Consistency, SDC），並結合潛在一致性匹配（Latent Consistency Matching, LCM），在不改變 MTP 模型架構的前提下，分別於輸出機率分布與潛在表示層級強化第一預測頭與其餘預測頭之間的一致性。SDC 以第一預測頭作為教師，利用其預測機率分布監督其餘學生預測頭，並透過信心度門檻排除低信心度位置，以降低不可靠教師訊號對一致性學習的影響；LCM 則透過均方誤差縮小教師與學生預測頭之隱藏表示差距。

    本研究以 OpenWebText 進行模型訓練，並透過各預測頭之 Oracle perplexity、平均可接受草稿詞元數（Average Draft Tokens Accepted）、生成文本 perplexity、entropy、4-gram repetition，以及跨資料集零樣本評估，分析所提方法對預測品質、草稿詞元接受能力與生成品質的影響。

    正式 1.3B 模型實驗結果顯示，相較於標準 MTP baseline，SDC 與 LCM 聯合訓練可降低四個預測頭的 perplexity，其中第四預測頭的降幅最大，達 5.9\%；平均可接受草稿詞元數提升 4.2\%；生成文本 perplexity 則降低 52.7\%。此外，跨資料集零樣本評估亦顯示，所提方法在多數資料集與預測頭上具有改善效果。整體而言，實驗結果顯示，透過輸出分布與潛在表示兩個層級的一致性學習，可改善較遠未來詞元預測頭的最佳化表現，並提升 MTP 的多詞元預測品質與草稿詞元接受能力，展現進一步提升解碼效率的潛力。

    % 這個Command會自動幫你把Config裡面設定的東東填進來
    \zhAbsKeywords
\end{zhAbstract}
